\input{../../../stock/en-tete_v6.tex}

\newcommand{\titrechapitre}{Vecteurs}
\newcommand{\numerochapitre}{3}

\renewcommand{\headrulewidth}{0.4pt}
\renewcommand{\footrulewidth}{0.4pt}
\fancyfoot[L]{Raphaël JONTEF}
\fancyfoot[C]{\thepage}
\fancyfoot[R]{\href{https://www.alveusclub.com/}{Alveus}}
\fancyhead[L]{Mathématiques : Seconde}
\fancyhead[R]{\titrechapitre}

\begin{document}
\input{../../../stock/commands.tex}

% Page de titre custom
\setlength{\headheight}{15pt}
\begin{titlepage}
    \centering
    \vspace*{3cm}
    {\huge Chapitre \numerochapitre\par}
    {\Huge \bfseries\titrechapitre\par}
    \vspace{8cm}
    {\Large Cours rédigé par Raphaël JONTEF\par}
    \vspace{2em}
    {\small \textit{raphaeljontef@hotmail.com}\par}
    \vspace{2cm}

    {\normalsize pour un enseignement dans le cadre  \par}
    {\small d'un soutien scolaire encadré par \href{https://www.alveusclub.com/}{Alveus} \par}
    \vfill
    {\large \today\par}
\end{titlepage}

% plan du cours
% \tableofcontents
\newpage


% -----------Corps du document--------------

\section{Vecteurs et géométrie du plan}

\subsection{Notion de vecteur}

\begin{definition}{}{vecteur}
Un \notion{vecteur} est défini par :
\begin{itemize}
    \item une \notion{direction} (celle de la droite qui le contient),
    \item un \notion{sens} (donné par la flèche),
    \item une \notion{longueur} (appelée \textit{norme}).
\end{itemize}
Le vecteur $\overrightarrow{AB}$ est défini par le déplacement du point $A$ vers le point $B$.
\end{definition}

\begin{exemple}{}{vecteurs égaux}
Les vecteurs $\overrightarrow{AB}$ et $\overrightarrow{CD}$ sont égaux si et seulement si ils ont :
\begin{itemize}
    \item la même direction,
    \item le même sens,
    \item et la même longueur.
\end{itemize}
\end{exemple}

\begin{remarque}{}{vecteur nul}
Le vecteur $\overrightarrow{AA}$ est appelé \notion{vecteur nul}, noté $\overrightarrow{0}$.  
Il n’a ni direction, ni sens, ni longueur.
\end{remarque}


\subsection{Somme et différence de vecteurs}

\begin{proposition}{}{relation de Chasles}
Pour tous points $A$, $B$ et $C$ du plan :
\[
\overrightarrow{AB} + \overrightarrow{BC} = \overrightarrow{AC}
\]
C’est la \notion{relation de Chasles}.
\end{proposition}

\begin{exemple}{}{relation de Chasles}
Si un mobile se déplace de $A$ vers $B$, puis de $B$ vers $C$, son déplacement total est le vecteur $\overrightarrow{AC}$.
\end{exemple}

\begin{proposition}{}{opposé d'un vecteur}
Deux vecteurs $\overrightarrow{u}$ et $\overrightarrow{v}$ sont \notion{opposés} si :
\[
\overrightarrow{u} + \overrightarrow{v} = \overrightarrow{0}
\]
Autrement dit, $\overrightarrow{BA} = -\overrightarrow{AB}$.
\end{proposition}

---

\subsection{Repérage dans le plan}

\begin{definition}{}{repère}
Un \notion{repère du plan} est défini par un point $O$ (l’origine) et deux vecteurs non colinéaires $\vec{i}$ et $\vec{j}$.  
Le repère est dit \notion{orthonormé} si $\vec{i}$ et $\vec{j}$ sont orthogonaux et de longueur $1$.
\end{definition}

\begin{definition}{}{coordonnées d'un vecteur}
Dans un repère $(O; \vec{i}, \vec{j})$ :
\[
\overrightarrow{AB} = (x_B - x_A \, ; \, y_B - y_A)
\]
\end{definition}

\begin{exemple}{}{coordonnées}
Soit $A(2; 1)$ et $B(5; 4)$.  
Alors :
\[
\overrightarrow{AB} = (5 - 2 \, ; \, 4 - 1) = (3; 3)
\]
\end{exemple}

\begin{proposition}{}{opérations vecteurs}
Si $\overrightarrow{u} = (x_1, y_1)$ et $\overrightarrow{v} = (x_2, y_2)$ :
\[
\overrightarrow{u} + \overrightarrow{v} = (x_1 + x_2, \, y_1 + y_2)
\quad \text{et} \quad
k\overrightarrow{u} = (kx_1, \, ky_1)
\]
\end{proposition}

\subsection{Produit scalaire}



\begin{definition}{}{produit scalaire}
    Le \notion{produit scalaire} de deux vecteurs $\overrightarrow{u} = (x_1, y_1)$ et $\overrightarrow{v} = (x_2, y_2)$ est défini par :
    \[
    \overrightarrow{u} \cdot \overrightarrow{v} = x_1 x_2 + y_1 y_2
    \]
\end{definition}
---

\subsection{Alignement, milieu, colinéarité}

\begin{definition}{}{colinearité}
Deux vecteurs $\overrightarrow{u}(x_1, y_1)$ et $\overrightarrow{v}(x_2, y_2)$ sont \notion{colinéaires} si :
\[
x_1 y_2 - x_2 y_1 = 0
\]
\end{definition}

\begin{proposition}{}{alignement}
Trois points $A$, $B$ et $C$ sont \notion{alignés} si et seulement si les vecteurs $\overrightarrow{AB}$ et $\overrightarrow{AC}$ sont colinéaires.
\end{proposition}

\begin{exemple}{}{alignement}
Soit $A(0;0)$, $B(1;2)$ et $C(2;4)$.  
On a $\overrightarrow{AB} = (1;2)$ et $\overrightarrow{AC} = (2;4)$, donc $2\times(1;2) = (2;4)$.  
Les vecteurs sont colinéaires : les points sont alignés.
\end{exemple}

\begin{proposition}{}{milieu}
Le milieu $M$ du segment $[AB]$ a pour coordonnées :
\[
M\left( \frac{x_A + x_B}{2} \, ; \, \frac{y_A + y_B}{2} \right)
\]
\end{proposition}

---

\subsection{Norme et distance}

\begin{definition}{}{norme}
La \notion{norme} d’un vecteur $\overrightarrow{u} = (x, y)$ est :
\[
\|\overrightarrow{u}\| = \sqrt{x^2 + y^2}
\]
\end{definition}

\begin{proposition}{}{distance}
La \notion{distance} entre deux points $A(x_A, y_A)$ et $B(x_B, y_B)$ est :
\[
AB = \sqrt{(x_B - x_A)^2 + (y_B - y_A)^2}
\]
\end{proposition}

\begin{exemple}{}{distance}
Si $A(1;2)$ et $B(4;6)$ :
\[
AB = \sqrt{(4 - 1)^2 + (6 - 2)^2} = \sqrt{9 + 16} = 5
\]
\end{exemple}

\begin{remarque}{}{Pythagore}
Cette formule est une application directe du théorème de Pythagore dans le triangle rectangle formé par les points $A$, $B$, et leur projection sur les axes.
\end{remarque}



\end{document}